\section{MobileGym++}
\label{sec:mobilegym++}
\label{sec:method}

\subsection{A gap in controllable commercial applications}
\label{sec:mgpp-env}

Mobile agents are increasingly moving toward real-world workflows across commercial applications, such as travel booking on Trip and online shopping on Taobao, where tasks are often long-horizon and grounded in real user needs. However, training and evaluation environments for such commercial applications remain largely missing, leaving agents with limited opportunities to practice and be evaluated in realistic mobile-use scenarios.

Real-device or Android-emulator setups face login, reset, and real-world side effects that make large-scale training impractical. Prior simulated-app work typically covers only a small slice of commercial apps and rarely provides resettable business state~\citep{rawles2025androidworld,kong2025mobileworld}.

We built MobileGym++, a browser-hosted mobile environment that carries commercial workflows in controllable simulated applications. The runtime extends MobileGym~\citep{wu2026mobilegym}, preserving its advantages in reset and programmatic evaluation while substantially expanding domain coverage, the number and richness of applications, and support for realistic cross-app workflows. Observations remain screen pixels, there is no accessibility tree, and the agent cannot read business state; device state can be exported and restored, and many instances can run in parallel. The current installable environment contains 64 ready-to-use applications spanning travel, shopping, finance, content, social, office, and local services; 35 of these applications were newly built in this round. Each application maintains persistent business states, such as orders, budgets, and itineraries, which evolve through interaction and enable task verification based on the resulting state.

Figure~\ref{fig:mgpp-overview} summarizes the environment at a glance: commercial-app coverage, executable single- and cross-app tasks, the two interaction channels, and their shared terminal check.

\begin{figure}[t]
\centering
\includegraphics[width=\linewidth]{figures/mobilegym_fig1_20260913.pdf}
\caption{MobileGym++ combines broad commercial-app coverage with executable single- and cross-app tasks and two interchangeable interaction channels. GUI-only and Hybrid share the same initial state and terminal check.}
\label{fig:mgpp-overview}
\end{figure}

Table~\ref{tab:mgpp-stack-compare} makes the comparison scope explicit: MobileGym++ is an environment with 64 installable applications, and MobileGym++ Bench is the 215-task stack that exercises 48 of them. The table uses common comparison axes: app and task scale, cross-app coverage, interaction modes, reproducibility, and deterministic evaluation. The last row reports the environment scale; the dagger records the 48-app Bench coverage.

\begin{table*}[t]
\caption{Comparison of mobile-agent benchmarks on scale, interaction, and evaluation. MobileGym++ uses a 64-application simulated environment of commercial-style apps; its 215-task Bench covers 48 of these applications. The Platform column distinguishes API, Android, and simulated settings. Following recent benchmark comparisons, capabilities are shown as explicit columns. $\checkmark$ denotes supported, $\times$ unsupported, and --- not reported or not applicable. The two MobileGym rows share the same reproducible simulation and deterministic state-based evaluation; scores are omitted because the platforms and task definitions do not share a protocol.}
\label{tab:mgpp-stack-compare}
\centering
\scriptsize
\setlength{\tabcolsep}{3.0pt}
\renewcommand{\arraystretch}{1.16}
\resizebox{\linewidth}{!}{%
\begin{tabular}{@{}l l r r c c c c c@{}}
\toprule
 & & \multicolumn{2}{c}{Scale} & \multicolumn{3}{c}{Interaction} & \multicolumn{2}{c}{Evaluation} \\
\cmidrule(lr){3-4} \cmidrule(lr){5-7} \cmidrule(lr){8-9}
Benchmark & Platform & \# Apps & \# Tasks & Cross-app & GUI & Native tools / MCP & Reproducible & Deterministic eval. \\
\midrule
AppWorld~\citep{trivedi2024appworld} & API & 9 & 750 / 250 sc. & $\checkmark$ & $\times$ & --- & $\checkmark$ & $\checkmark$ \\
AndroidLab~\citep{xu2025androidlab} & Android & 9 & 138 & --- & $\checkmark$ & --- & --- & $\checkmark$ \\
AndroidWorld~\citep{rawles2025androidworld} & Android & $\sim$20 & 116 tmpl. & $\checkmark$ & $\checkmark$ & --- & $\checkmark$ & $\checkmark$ \\
SPA-Bench~\citep{chen2024spabench} & Android & 66 & 340 & $\checkmark$ & $\checkmark$ & --- & $\checkmark$ & $\times$ \\
MobileWorld~\citep{kong2025mobileworld} & Android & 20 & 201 & $\checkmark$ & $\checkmark$ & $\checkmark$ & $\checkmark$ & $\checkmark$ \\
MobileGym-Bench~\citep{wu2026mobilegym} & Sim. & 28 & 416 tmpl. & $\checkmark$ & $\checkmark$ & --- & $\checkmark$ & $\checkmark$ \\
\textbf{MobileGym++ Bench (ours)} & \textbf{Sim.} & \textbf{64$^\dagger$} & \textbf{215} & \textbf{$\checkmark$} & \textbf{$\checkmark$} & \textbf{$\checkmark$} & \textbf{$\checkmark$} & \textbf{$\checkmark$} \\
\bottomrule
\end{tabular}%
}
\vspace{2pt}
\begin{flushleft}
\footnotesize
$^\dagger$ The environment contains 64 installable applications; the frozen Bench exercises 48 of them. --- means that the cited benchmark does not report the capability in a comparable form. Deterministic evaluation refers to an automatic, reproducible success check, not a model judge.
\end{flushleft}
\end{table*}

\subsection{35 high-fidelity commercial applications}
\label{sec:mgpp-apps}

The 35 applications added above were implemented from the corresponding real apps through a human-annotator and coding-agent pipeline: annotators collect screen recordings and long screenshots; the agent reproduces the client; humans check the result against the real device and send revisions back through the same tools. Details are in Appendix~\ref{app:mgpp-build}.

The newly built applications have sufficient richness and interaction depth. On these clients, the median navigation declaration is 37 routes, 175 interactive actions, and 88 transitions; median interaction depth from the root is 3. As a depth reference, the thickest application in MobileGym has 100 routes and 265 actions, and its deepest terminal takes 7 steps to reach. The function channel is described in Section~\ref{sec:mgpp-tools}.





\subsection{Realistic cross-application workflows}
\label{sec:mgpp-crossapp}

Several commercial applications sit on one resettable device, each maintaining its own live business state: goods and orders in shopping, expenses and budgets in bookkeeping, itineraries and changes in the calendar, contacts and addresses in the address book. These objects change with on-screen actions and leave a trace in device state.

Clipboard, system broadcasts, and inter-app navigation connect the apps. A workflow can leave the current app, enter a second app, and rewrite that app's own business objects: a payment in shopping can produce a matching expense in bookkeeping; a selected train number in the calendar can enter an itinerary app for a change; an address from contacts can enter a checkout flow. The second app therefore receives a real business write, not only a forwarded message.

Every step of a cross-app workflow acts on the same resettable device state, so the same workflow can be replayed.

\subsection{Hybrid GUI and app-native tool use}
\label{sec:mgpp-tools}

GUI-only is already a complete setting. To support longer-horizon study we additionally open an app-native function channel on the same application state, so when to use the interface, when to call a function, and how to return to the interface for confirmation can be examined in one world. Hybrid is another completion path, not a requirement.

Approximately 52 applications expose a function interface in the foreground, occupying the same action slot as a tap, swipe, or text input. Functions are a subset of GUI-reachable operations: prefill creates a draft only, while saving, payment confirmation, passwords, and contact selection return to the GUI. Tools are used mainly to locate and organize structured information; both channels write the same state.

\begin{figure}[t]
\centering
\includegraphics[width=\linewidth]{figures/mobilegym_fig3_20260913.pdf}
\caption{Hybrid execution in MobileGym++. GUI actions and app-native tools operate on the same business state; a trajectory may hand off between channels before the checker reads the resulting state diff.}
\label{fig:mgpp-hybrid}
\end{figure}

\subsection{MobileGym++ Bench}
\label{sec:mgpp-bench}

Existing mobile-agent benchmarks cover many single-app GUI operations and increasingly include cross-application workflows and tool interaction~\citep{rawles2025androidworld,kong2025mobileworld}. MobileGym++ Bench asks whether an agent can finish long-horizon, realistic, cross-application commercial work in a controllable commercial environment. Task length comes from locating, filtering, confirming, paying, and writing back across apps.



\paragraph{Task design.}
Tasks are organized around executable user goals and checkable terminals. Core tasks close with a business operation such as placing an order, posting an expense, creating an itinerary, sending a message, or modifying an app object. The result of that operation lands in the application's own business state. A single-app task reaches that write only after locating and filling fields across deeper pages. Instructions use user-goal language and do not name widget labels, tool namespaces, or checker fields.

In a cross-application task, the value needed for a write is present in another app's live state rather than written as a fixed constant in the instruction. The agent must identify the source, carry the value across the app boundary, and write it into the correct field; the checker verifies the relation between the resulting business fields.

\paragraph{Scale and composition.}
The current set contains 215 tasks across seven business domains: 109 mention one app, 74 two, 27 three, and 5 four. Tasks are distributed over multiple applications and business objects, so evaluation is not dominated by a handful of apps.

\paragraph{Evaluation protocol.}
Each task has three parts: a natural-language user goal, a bound initial state, and a programmatic terminal check. Success is decided by a programmatic checker: it reads the before/after diff of application business state, does not read the trace or screenshots, and does not use a model judge; different routes that accomplish the same write are equivalent. Allowed state paths are declared, and other mutations are recorded separately from success. The primary metric is task success rate. GUI-only and Hybrid are evaluated on the same tasks, initial states, and terminal checks. Mean steps on successful trajectories measure interaction efficiency conditional on completion; because failed trajectories often reach the step cap, we also report the success rate and the fraction reaching the cap to separate efficiency from reliability.
